

% \section{Preliminaries: Standard Reinforcement Learning}
% \label{subsec:preliminaries}

% Reinforcement Learning aims to optimize a policy $\pi_{\theta}$ to maximize the expected cumulative reward $J(\theta) = \mathbb{E}_{\tau \sim \pi_{\theta}}[R(\tau)]$, where $R(\tau)$ is the cumulative reward of trajectory $\tau$. A standard approach to optimize this objective is the Policy Gradient method \citep{reinforce, policy-gradient}, which updates parameters in the direction of the gradient:

% \begin{equation}
%     \nabla_{\theta} J(\theta) = \mathbb{E}_{s,a \sim \pi_{\theta}} [\nabla_{\theta} \log \pi_{\theta}(a|s) \cdot A(s,a)]
%     \label{eq:policy_gradient}
% \end{equation}

% where $A(a,s)$ is the advantage function, formulated as: 

% \begin{equation}
%     A(s,a) = Q(s,a) - V(s),
%     \label{eq:advantage_def}
% \end{equation}

% quantifying under the state $s$, how much better the performance of a specific action $Q(s,a)$ is compared to the average performance $V(s)$. 

% In standard Actor-Critic architectures, obtaining $Q(s,a)$ and $V(s)$ typically relies on training parameterized value networks, which is computationally expensive and yields a static model.


\section{Preliminaries}
\label{subsec:preliminaries}

Reinforcement Learning (RL) aims to optimize a policy $\pi_{\theta}$, which maps a given state to a probability distribution over actions, to maximize the expected cumulative reward. The objective function is defined as $J(\theta) = \mathbb{E}_{\tau \sim \pi_{\theta}}[R(\tau)]$, where $\tau = (s_0, a_0, r_0, s_1,a_1,r_1\dots, s_n, a_n,r_n)$ represents a trajectory of $n$ interactions, and $R(\tau)$ denotes the cumulative discounted return of that trajectory. 

A fundamental approach to optimize this objective is the Policy Gradient method \citep{reinforce, policy-gradient}. Intuitively, this method shifts the probability mass towards actions that yield higher returns by updating parameters in the direction of the gradient:
\begin{equation}
    \nabla_{\theta} J(\theta) = \mathbb{E}_{s,a \sim \pi_{\theta}} [\nabla_{\theta} \log \pi_{\theta}(a|s) \cdot A(s,a)],
    \label{eq:policy_gradient}
\end{equation}
where $A(s,a)$ is the \textbf{advantage function}. The advantage function serves as a critic to evaluate the selected action relative to a baseline, formulated as:
\begin{equation}
    A(s,a) = Q(s,a) - V(s).
    \label{eq:advantage_def}
\end{equation}
Here, $Q(s,a)$ is the \textbf{action-value function}, representing the expected return of taking action $a$ in state $s$, while $V(s)$ is the \textbf{state-value function}, representing the average expected return of being in state $s$. Thus, $A(s,a)$ quantifies \textit{how much better} a specific action is compared to the average performance of the policy in that state.

Reinforcement Learning typically requires the training of additional value networks to estimate the advantage. This process is computationally expensive and results in a static model that lacks the flexibility to adapt at test time.